\scanpage{13}
\mspage{27}

garantit alors que $v$ est un isomorphisme d'algèbres topologiques (grâce
à \unsure{c)} et d) pour $B$ et $C$). Donc \deleted{les} on a 2), \deleted{et}
\inserted{donc on a 3). Or} le même argument en
sens inverse montre que \deleted{$\unsure{u:C\longrightarrow B}$ est alors un isomorphisme d'alg top, i.e.
qu'on a 4)} \inserted{3) implique 4). -- Enfin $4)\Longrightarrow5)$ et $5)\Longrightarrow2)$.}
On peut certainement améliorer la rédaction de la démonstration, évitant certaines redites.

Pour simplifier, nous nous restreindrons par la suite à des hypothèses
noethériennes (permettant de nous débarrasser en particulier de la désagréable hypothèse d) sur $B$).

\uline{Définition} \inserted{4.}7. Soit $A\longrightarrow B$ un homomorphisme continu d'anneaux noethériens
préadiques. On dit que $B$ est une \uline{algèbre} (topologique) \uline{de Cohen} sur $A$
si \inserted{a)} $B$ est séparé et complet, i.e. \inserted{est} un anneau adique,
\deleted{\unsure{pourvu}} \inserted{b)} s'il satisfait à la condition P, et si enfin \inserted{c)} $B$ est formellement
\deleted{\unsure{plat}} \inserted{simple} sur $A$, \deleted{\unsure{ou}}
\inserted{(ou encore, $I$ étant un idéal de définition de $A$, $B_0=B/IB$ est formel\unsure{t} simple sur $A_0=A/IA$)}

On notera qu'en vertu de 3.1., si $I$ est un idéal de définition de $A$,
la dernière condition équivaut à la condition que $B/IB$ est formellement
simple sur $A/IA$. D'ailleurs, on peut interpréter \inserted{4.}6. de la façon
suivante (et il est \deleted{peut-être} plus raisonnable de reléguer \inserted{4.}6. en lemme,
et de réserver le corollaire en théorème):

\uline{Corollaire} \inserted{4.}8. Soient $A$ un anneau \deleted{préadique} noethérien, $I$ un idéal de $A$,
munissons $A$ de la topologie $I$-préadique, posons $A_0=A/IA$, et soit $B_0$ une
$A_0$-algèbre topologique. \deleted{Si il}\inserted{Supposons qu'il} existe une algèbre de Cohen $B$ sur $A$ telle
que $B/IB$ soit isomorphe à $B_0$ comme $A_0$-algèbre topologique,
\inserted{(ce qui implique $B_0$ \deleted{\unsure{formellement simple}}\inserted{\unsure{anneau de Cohen}} sur $A_0$).}
alors $B$ est unique à \inserted{un} isomorphisme près
\deleted{(induisant l'isomorphisme donné sur} (compatible avec les isomorphismes donnés des algèbres
\deleted{réduites} réduites mod $I$ avec $B_0$). De plus, si $B$ est une $A$-algèbre topologique adique telle que $B/IB$
soit isomorphisme à $B_0$ comme $A_0$-algèbre topologique, alors les conditions
suivantes sur $B$ sont équivalentes: \inserted{a)} $B$ est une $A$-algèbre de Cohen

\mspage{12}
